\section{The CASIA v2.0 Dataset and Its Pitfalls}\label{sec:dataset}

We use CASIA~v2.0~\cite{dong2013casia}, which has 12,614 images: 7,491 authentic and 5,123 tampered. The tampered
images were made with editing software as either \emph{copy-move} (a region duplicated within one image) or
\emph{splicing} (a region pasted from a different image). There are nine scene categories, and most images are
small: 85.8\% of authentic and 64.7\% of tampered images are $384\times256$ or $256\times384$ pixels.

\subsection{Labels and Masks}\label{sec:dataset-labels}

The forgery type is normally read from the file-name letter (Tp\_S = copy-move, Tp\_D = splicing), which gives
3,274 copy-move and 1,849 splicing images. For 115 images this letter conflicts with other evidence: the letter in
the mask's file name, and whether the host and donor ids are equal (equal ids mean copy-move). We take the majority
of these three votes, which changes 99 labels (60 splicing $\rightarrow$ copy-move, 39 copy-move $\rightarrow$
splicing). The corrected type (Table~\ref{tab:dataset-format}) is used for stratification and all per-type
results.

\begin{table}[t]
\centering
\caption{Class and file format (corrected forgery type)}\label{tab:dataset-format}
\footnotesize
\begin{tabular}{lrrrr}
\toprule
Class & Images & JPEG & TIFF & BMP \\
\midrule
Authentic & 7,491 & 7,437 & 0 & 54 \\
Copy-move & 3,295 & 967 & 2,328 & 0 \\
Splicing & 1,828 & 1,097 & 731 & 0 \\
\bottomrule
\end{tabular}
\end{table}

Each tampered image has a binary mask marking the content taken from the donor. Because 142 mask names differ
from their image name, we match masks on the trailing tamper id, which is unique across the dataset; every tampered
image then has exactly one mask. Of the 5,123 masks, 19 are unusable for pixel-level scoring (5 transposed, 12
with other sizes, 2 covering more than 95\% of the image). These images remain in the classification experiments
but are excluded from localisation scoring, which leaves 5,104 valid masks. Checked against where the tampered image
differs from its background image, the 5,034 checkable masks have median recall 0.885 and median precision 0.991,
so the masks are generally accurate; 52 masks with recall $<0.2$ are flagged but not removed.

Tampered regions are small: the median tampered area is 3.1\% of the image (10th percentile 0.5\%, 90th
percentile 28\%). We therefore report localisation in three area buckets: $<1\%$ (1,114 images), 1--5\%
(2,023) and $>5\%$ (1,967). A tampered file name names two authentic images; we treat the first-named one as the
\emph{host}, since it supplies most pixels in 1,689 of the 1,778 splicing-named images with differing ids. In 23
images the second-named image supplies most pixels (18 large pastes and 5 apparently swapped ids); these are
additionally grouped with the second-named image (Section~\ref{sec:leakage}). Details of the mask audit and the
host analysis are in Appendix~\ref{app:dataset}.

\subsection{Pitfalls That Affect Evaluation}\label{sec:dataset-pitfalls}

\emph{(1) File format tracks the label.} 60\% of tampered images are uncompressed TIFF, but no authentic image is
(Table~\ref{tab:dataset-format}). \emph{(2) JPEG settings track the label.} Among JPEG files, 89.1\% of authentic
images use a standard libjpeg (IJG) quantisation table, compared with 0\% of tampered images, whose tables look like
those written by editing software. The closest IJG quality factor is mostly 84 or 90 for authentic images and 93
for tampered ones. \emph{(3) Image size tracks the label.} Tampered images are more often $640\times480$ or
$800\times600$. \emph{(4) Duplicates.} 196 authentic images appear twice under different ids with identical
pixels, and about 480 tampered images are near-identical to an authentic image other than their named host.
\emph{(5) Label noise in names}, corrected as described above.

Points 1--3 mean that a classifier can separate the classes without looking at any forensic evidence.
Section~\ref{sec:leakage} measures this directly.
